% ===================================================================
%  CHART No. 1 — The School. — The Secondary School. — The Classroom.
%
%  Ceci n'est PAS une transcription : c'est une traduction, faite en
%  2026, en anglais d'aujourd'hui. D'ou trois differences avec
%  texto/io et texto/fr, et trois seulement :
%
%   * pas de \nl ni de \cc — la traduction ne suit aucune ligne
%     imprimee, et les lignes de ce fichier ne sont que du confort ;
%   * pas de \begin{VUpage} — elle n'a ni feuillet ni folio, et la
%     page de lecture ne lui en affiche pas ;
%   * le gras se pose sur le TERME seul, ponctuation dehors.
%
%  Tout le reste est identique, et doit l'etre : les cles « %%K » sont
%  celles de texto/io, macro pour macro pour l'apparat, et les renvois
%  \textsuperscript{(n)} visent les MEMES objets de la planche — ce
%  sont eux qui ouvrent les gros plans.
%
%  L'ido est le texte de base : c'est lui que la colonne de gauche
%  porte, et c'est donc lui qu'on traduit. Le francais sert de
%  controle — les deux editions disent la meme chose a quelques mots
%  pres — et il tranche quand l'ido est elliptique. La ou l'ido
%  subdivise et pas le francais (« La bona lernanti », « La mala
%  lernanti »), l'anglais suit l'ido.
% ===================================================================

%%K t01-apar-0 apar
\VUsaut{2.0mm}
\VUcentre{12.6pt}{120}{EXPLANATORY BOOKLET}
\VUsaut{3.2mm}
\VUcentre{8.4pt}{160}{FOR}
\VUsaut{2.6mm}
\VUtitre{15.6pt}{0}{the Delmas Auxiliary Charts}
\VUsaut{3.4mm}
\VUornamento{9pt}
\VUsaut{5.0mm}
\VUcentre{13.2pt}{90}{FIRST SERIES}
\VUsaut{3.0mm}
\VUfilet{16mm}
\VUsaut{5.6mm}

%%K t01-apar-1 apar
\VUtitre{15.0pt}{60}{CHART No. 1}
\VUsaut{1.4mm}
\VUtitre{11.4pt}{0}{The School. --- The Secondary School. --- The Classroom.}
\VUsaut{2.2mm}
\VUfilet{20mm}
\VUsaut{6.4mm}

%%K t01-tit-1 sub
\VUpk{10.2pt}{40}{General description.}
\VUsaut{3.0mm}

%%K t01-01-1 p
1. --- This chart shows a \VUgras{classroom} in a \VUgras{secondary school}.
In this classroom there are eight \VUgras{tables} \textsuperscript{(18)} and eight
\VUgras{benches} \textsuperscript{(32)}. Three pupils sit on each bench. The
\VUgras{teacher} \textsuperscript{(7)} sits on a \VUgras{chair} \textsuperscript{(8)} at his
\VUgras{desk} \textsuperscript{(6)}.

%%K t01-02-1 p
2. --- To the left of the desk stands a glass-fronted \VUgras{cabinet}
\textsuperscript{(15)}. To the right is the \VUgras{blackboard} \textsuperscript{(1)} with the
\VUgras{cloth} \textsuperscript{(2)} and the \VUgras{chalk} \textsuperscript{(3)}.

%%K t01-02-2 p
Above the desk hang \VUgras{wall charts} \textsuperscript{(9, 11, 12)} and a
\VUgras{map} \textsuperscript{(10)}.

%%K t01-03-1 p
3. --- Not all the pupils are in their \VUgras{seats} \textsuperscript{(17)}.
John \textsuperscript{(4)} is standing at the board; he is holding the cloth and
wiping off the \VUgras{geometrical figures}.

%%K t01-03-2 p
Peter is near the desk; he is collecting the \VUgras{written exercises}
\textsuperscript{(5)} and bringing them to the teacher.

%%K t01-04-1 p
4. --- \VUgras{This} teacher is the teacher of \VUgras{modern languages}. He
teaches the pupils to speak, read and write foreign languages
\VUgras{(German, English, Spanish, Italian, Russian} or \VUgras{French)}.

%%K t01-05-1 p
5. --- The classroom is very large and very bright. At the far end
there is a wide \VUgras{window} with two \VUgras{transoms} \textsuperscript{(78)} and a
\VUgras{glazed door} to air it and light it.

%%K t01-05-2 p
This classroom is not cold. In the middle, to heat it, there is a very
good \VUgras{stove} \textsuperscript{(46)} with its \VUgras{pipe} \textsuperscript{(45)}.

%%K t01-05-3 p
\VUgras{Coat hooks} \textsuperscript{(58)} are fixed to the partition; the pupils' hats
and coats hang from them.

%%K t01-tit-2 sub
\VUsaut{5.2mm}
\VUpk{10.2pt}{40}{The playground.}
\VUsaut{3.6mm}

%%K t01-06-1 p
6. --- The \VUgras{clock} \textsuperscript{(63)} says five past nine.

%%K t01-06-2 p
\VUgras{Break} \textsuperscript{(76)} has just ended and the lesson is starting again.
Break is held in the \VUgras{yard} \textsuperscript{(75)}. We can see a few pupils
playing. An \VUgras{assistant master} \textsuperscript{(74)} is supervising them.

%%K t01-07-1 p
7. --- In the yard we can also see various other people: the
\VUgras{inspector} \textsuperscript{(73)}, the \VUgras{headmaster} \textsuperscript{(71)} and the
\VUgras{deputy head} \textsuperscript{(72)}.

%%K t01-07-2 p
The inspector inspects the schools, the headmaster runs the school, the
deputy head oversees the pupils' work.

%%K t01-07-3 p
The fourth figure is the \VUgras{caretaker} \textsuperscript{(77)}. He is carrying a
register under his arm to record the absentees.

%%K t01-08-1 p
8. --- At the far end of the yard are the \VUgras{gymnastic apparatus}
\textsuperscript{(70)} and the workshop for \VUgras{handicrafts} \textsuperscript{(69)}.

%%K t01-08-2 p
To the right of the chart we can just make out the \VUgras{classrooms} for
\VUgras{music} and \VUgras{drawing}.

%%K t01-noto-1 noto
\VUnotes{143.2mm}{%
(*) For \textit{first names} it is also recommended to transcribe them in
their Latin form. That is the only way to avoid countless variants.
Besides, how would one choose between \textit{Giovanni, Jean, Johann,}
\textit{Jan, Hans, John, Ivan,} and so on? Are we to create a special
dictionary of first names? Let us take the Latin nominative and say:
\VUgras{Ioannes, Ioanna; Iulius, Iulia; Iakobus; Andreas; Lukas.}
(Complete Grammar).%
}

%%K t01-tit-3 sub
\VUpk{10.2pt}{40}{Detailed description.}
\VUsaut{2.4mm}

%%K t01-tit-4 sub
\VUpk{10.2pt}{40}{The good pupils.}
\VUsaut{3.0mm}

%%K t01-09-1 p
9. --- The pupils on the first bench to the right of the desk are
\VUgras{Henry} (*), \VUgras{Stephen} and \VUgras{Charles}.

%%K t01-09-2 p
Henry is very attentive and hard-working; he is listening to the
teacher. On his table he has an \VUgras{exercise book} \textsuperscript{(28)}, a
\VUgras{book} \textsuperscript{(29)}, a \VUgras{dictionary} \textsuperscript{(30)} and a
\VUgras{pencil case} \textsuperscript{(31)}. His book has many \VUgras{pages}; each chapter
contains several \VUgras{paragraphs} and each \VUgras{line} many \VUgras{words}.

%%K t01-09-4 p
His neighbour Stephen \textsuperscript{(24)} is keen. He is writing down in his
exercise book the lesson set for next time. He has fine
\VUgras{handwriting}. His book is closed. We can see only the \VUgras{cover}
\textsuperscript{(27)}. Beside him lies his pair of \VUgras{compasses} \textsuperscript{(26)}.

%%K t01-09-5 p
Charles \textsuperscript{(23)} is holding his book in his hand; he is opening it to
go over his lesson again. Like every pupil, he has an \VUgras{inkwell}
\textsuperscript{(25)} in front of him. He is obedient.

%%K t01-10-1 p
10. --- At the next table are \VUgras{Louis, Anthony} and \VUgras{Paul}. Louis
is reciting his \VUgras{lesson} \textsuperscript{(22)} and answering the teacher's
\VUgras{questions}. Anthony \textsuperscript{(21)} is very inattentive; sometimes the
teacher even punishes him. Paul \textsuperscript{(20)} is a good friend, pleasant and
obliging. He is very attentive.

%%K t01-tit-5 sub
\VUsaut{4.2mm}
\VUpk{10.2pt}{40}{The bad pupils.}
\VUsaut{3.2mm}

%%K t01-11-1 p
11. --- Behind Henry sits \VUgras{George} \textsuperscript{(38)}. He is not a bad lad,
but he is \VUgras{talkative}, inattentive and fidgety. His \VUgras{frivolity}
and \VUgras{disobedience} cause \VUgras{disruption} during the lesson, and he
often deserves a stern \VUgras{warning}. His \VUgras{rough book}
\textsuperscript{(35)} is dirty, he covers it with \VUgras{ink blots} from his \VUgras{pen}
\textsuperscript{(34)}, which is why he always uses his \VUgras{eraser} \textsuperscript{(36)}; his
\VUgras{satchel} \textsuperscript{(33)} is torn, for he is very careless. Why does he
bring his \VUgras{pencil holder} \textsuperscript{(37)} to the modern languages
classroom?

%%K t01-11-2 p
His neighbour Edmund \textsuperscript{(39)} does not like him. He is admittedly a
little lazy, but he is quiet and still. He does not chatter; only,
instead of listening, he prefers to read the \VUgras{stories} in his
\VUgras{reader}.

%%K t01-11-3 p
The third pupil on this bench is \VUgras{Ludovic} \textsuperscript{(40)}. He is diligent.
He is writing on a white \VUgras{sheet of paper}. In front of him he has
laid his \VUgras{blotting paper} \textsuperscript{(41)}.

%%K t01-12-1 p
12. --- The pupils on the second bench, to the teacher's left, are very
young. The first, little \VUgras{Emil}, cannot write well yet; he brings his
\VUgras{slate} \textsuperscript{(42)}, his \VUgras{slate pencil} and his \VUgras{sponge}
\textsuperscript{(43)}.

%%K t01-12-2 p
Next to him are \VUgras{Augustus} and \VUgras{Bernard} \textsuperscript{(44)}. The latter
works well, but his \VUgras{appearance} is very untidy.

%%K t01-13-1 p
13. --- Behind them sit three \VUgras{idlers}. \VUgras{Albert} does not learn
his lessons. \VUgras{James} \textsuperscript{(47)} is asleep instead of listening. His
\VUgras{idleness} earns him \VUgras{reprimands} and even \VUgras{punishments}.
Lastly \VUgras{Christian} \textsuperscript{(48)} is far too flippant. He \VUgras{behaves}
badly and makes no effort to mend his ways.

%%K t01-14-1 p
14. --- Their neighbours, by contrast, are well behaved. \VUgras{Ernest}
\textsuperscript{(51)} likes order and regularity; he always uses his \VUgras{ruler}
\textsuperscript{(50)} and his \VUgras{pencil} \textsuperscript{(49)}. He is a \VUgras{day pupil}.

%%K t01-14-2 p
\VUgras{Francis} \textsuperscript{(52)} is a \VUgras{boarder}; he wears the uniform, and eats
and sleeps at the school. He is a good \VUgras{friend}.

%%K t01-14-3 p
Maurice is sharpening his \VUgras{pencil} with his \VUgras{penknife}
\textsuperscript{(56)}; he is very careful.

%%K t01-14-4 p
He brings all his books, his \VUgras{grammar} \textsuperscript{(55)}, his
\VUgras{notebook} \textsuperscript{(54)}, and even a \VUgras{writing pad} \textsuperscript{(53)}. His
\VUgras{school bag} \textsuperscript{(57)} is on the floor behind him.

%%K t01-14-5 p
The two tables at the back of the classroom are taken by
\VUgras{good pupils} \textsuperscript{(19)}.

%%K t01-tit-6 sub
\VUsaut{4.6mm}
\VUpk{10.2pt}{40}{Drawing and music.}
\VUsaut{3.4mm}

%%K t01-15-1 p
15. --- In the \VUgras{drawing room} there are handsome \VUgras{models} hung on
the wall and a \VUgras{lamp} \textsuperscript{(62)} with a \VUgras{shade}, hanging from the
ceiling. The \VUgras{teacher} \textsuperscript{(60)} is very patient; he corrects the
pupils' \VUgras{drawings} \textsuperscript{(61)} and teaches them to draw a \VUgras{bust}
\textsuperscript{(59)} from life.

%%K t01-16-1 p
16. --- The \VUgras{music room} looks very interesting. There one learns to
read \VUgras{music}, to sing the \VUgras{scale notes} \textsuperscript{(67)} written on the
\VUgras{stave} (do, re, mi, fa, sol, la, si\,=\,C. D. E. F. G. A. B.), to
know the \VUgras{clefs} and to sing charming \VUgras{melodies}. The scores are
placed on the \VUgras{music stand} \textsuperscript{(65)}. One of the pupils is
\VUgras{playing the piano} \textsuperscript{(64)} and the \VUgras{music teacher}
\textsuperscript{(66)} is \VUgras{beating time} \textsuperscript{(68)}.

%%K t01-17-1 p
17. --- We can just see a corner of the \VUgras{workshop} for
\VUgras{handicrafts} \textsuperscript{(69)}, where the children may learn to plane wood
and file iron. Not every school has one.

%%K t01-tit-7 sub
\VUsaut{5.0mm}
\VUpk{10.2pt}{40}{The science room.}
\VUsaut{3.6mm}

%%K t01-18-1 p
18. --- The mathematics teacher teaches the pupils \VUgras{geometry,}
\VUgras{arithmetic, calculation} and the \VUgras{metric system}. On the chart we
can see the main geometrical figures: a \VUgras{circle} \textsuperscript{(a)}, a
\VUgras{square} \textsuperscript{(b)}, a \VUgras{triangle} \textsuperscript{(c)}, a \VUgras{line}
\textsuperscript{(d)}.

%%K t01-18-2 p
We can also see a \VUgras{sum}; it is not an \VUgras{addition}, nor a
\VUgras{subtraction}, nor a \VUgras{multiplication}: it is a \VUgras{division}
\textsuperscript{(e)}.

%%K t01-19-1 p
19. --- On the metric system chart we can see: a \VUgras{metre}
\textsuperscript{(a)}, a \VUgras{balance} \textsuperscript{(b)} and its \VUgras{weights}, a
\VUgras{litre} \textsuperscript{(e)}, a \VUgras{decalitre} \textsuperscript{(f)}, a \VUgras{decilitre}
and a \VUgras{cubic metre} \textsuperscript{(d)}.

%%K t01-19-2 p
The pupils also study \VUgras{natural science} \textsuperscript{(11)} --- zoology
\textsuperscript{(a)}, botany \textsuperscript{(b)}, geology \textsuperscript{(c)} --- and \VUgras{history}
\textsuperscript{(12)}.

%%K t01-19-3 p
In the cabinet are the instruments for \VUgras{physics} \textsuperscript{(15)} and for
\VUgras{chemistry} \textsuperscript{(16)}.

%%K t01-19-4 p
On top of the cabinet are: a \VUgras{sphere} \textsuperscript{(14)}, a \VUgras{cone} and a
\VUgras{globe} \textsuperscript{(13)}.

%%K t01-19-5 p
Above the charts hangs a \VUgras{map} showing the five parts of the world:
\VUgras{Europe} \textsuperscript{(g)}, \VUgras{Asia} \textsuperscript{(e)}, \VUgras{Africa}
\textsuperscript{(f)}, \VUgras{America} \textsuperscript{(ab)} and \VUgras{Oceania} \textsuperscript{(h)}, and
the two great oceans, the \VUgras{Pacific} \textsuperscript{(c)} and the \VUgras{Atlantic}
\textsuperscript{(d)}.
\VUsaut{6.0mm}
\VUfilet{22mm}
